\item \textbf{Queratómetro para medición de la córnea.}
Está entrenando para convertirse en asistente de un técnico óptico. Un día, está aprendiendo cómo colocar una lente de contacto en el ojo de un paciente. Realiza una medición con un queratómetro, que mide la curvatura frontal de la córnea. Este instrumento coloca un objeto iluminado de tamaño conocido a una distancia conocida $p$ de la córnea. La córnea refleja la luz, actuando como un espejo, y la magnificación lateral $M$ de la imagen se mide mediante un telescopio. Determine el radio de curvatura $R$ de la córnea para las mediciones que realiza para el paciente: $p = 30.0\text{ cm}$ y $M = 0.0130$.